\chapter{Hypothesis 3: Robustness to Heavy-Tailed Noise}
\label{ch:hyp3}

\Cref{ch:hyp1,ch:hyp2} established that the phase-plane estimator recovers
sub-wavelength displacement and material change on clean, synthetic data.
This chapter asks whether that still holds once the data are corrupted by
realistic, heavy-tailed noise, and --- critically --- shows that the answer
depends on \emph{which} migration technique is used.

\paragraph{Hypothesis 3.} Multi-dimensional phase-plane regression can deal
with noisy data, given that the right migration technique is used.

\section{A Laplace Noise Model from Real Field Data}
\label{sec:hyp3-laplace}

Rather than injecting arbitrary synthetic noise, the noise level and shape
used throughout this thesis are fitted to real GPR field noise. A sample of
field noise is tracked through the processing pipeline and fitted with both
a Laplace and a Gaussian distribution at two stages: after the final
processing stage ("9 Crop Samples"), and at the last stage before the
spherical-gain correction ("7 Constant Velocity"), since the gain correction
inflates the amplitude scale by several orders of magnitude and is not
representative of the raw simulated $E_z$ amplitudes used elsewhere in this
thesis.

\begin{table}[htbp]
  \centering
  \caption{Laplace and Gaussian fits to real field noise, at two pipeline
    stages ($n=976\,244$ samples each). The heavier-tailed Laplace
    distribution is adopted throughout this thesis; loc $=0$ for both
    stages.}
  \label{tab:laplace-fit}
  \begin{tabular}{lcc}
    \toprule
    Stage & Laplace scale & Gaussian $\sigma$ \\
    \midrule
    9 Crop Samples (post-gain)        & $117\,584.0$ & $171\,156.1$ \\
    7 Constant Velocity (pre-gain)    & $6.085$      & $9.233$ \\
    \bottomrule
  \end{tabular}
\end{table}

In both cases the fitted distribution is heavier-tailed than a Gaussian of
matched variance, consistent with field GPR noise being dominated by
occasional large-amplitude clutter and interference rather than purely
thermal noise. The pre-gain fit is the one actually used to generate
synthetic noise for the experiments in \cref{sec:hyp3-lateral,sec:hyp3-vertical,sec:hyp3-diagonal,sec:hyp3-fluidflow}:
its Laplace \emph{shape} (loc $=0$, heavier tails than Gaussian) is kept,
but its scale is rescaled so that the resulting noise standard deviation is
exactly $10\,\%$ of each synthetic B-scan's own signal standard deviation
--- a light, realistic noise level rather than the raw fitted scale, which
would be incommensurate with the synthetic $E_z$ amplitudes.

\section{Migration-Algorithm Response to Pure Noise}
\label{sec:hyp3-purenoise}

Before testing noise robustness on real signal-plus-noise data, a sanity
check establishes what each migration algorithm does to noise \emph{alone}:
an empty B-scan (no scatterer, no real signal) containing only the Laplace
noise of \cref{sec:hyp3-laplace} is migrated with Kirchhoff, Gazdag, and
back-propagation. If a method turns pure noise into something that looks
like a coherent, scatterer-like focus, every noisy result elsewhere in this
thesis carries a false-positive risk that must be accounted for.

\textbf{Result: the three methods do not fail the same way.} Kirchhoff
turns pure noise into clearly coherent, smooth wave-like bands that could
easily be misread as real layered structure: its delay-and-sum aperture
stacking imposes coherence on incoherent input by construction, summing
many traces along travel-time curves and smoothing incoherent noise into
locally-correlated structure. Gazdag's noise output, by contrast, stays
speckled and incoherent, with no wave-like artefacts --- just texture. This
is a direct, practical consequence of the structural difference between the
two algorithms (\cref{sec:th-migration}): Kirchhoff's spatial stacking
manufactures apparent coherence from nothing, while Gazdag's
frequency-domain downward continuation does not.

\draftnote{the pure-noise Kirchhoff/Gazdag comparison above is implemented
and run in \texttt{Noise\_Playground.ipynb}, but its plots have not yet been
exported to a \texttt{TimeLapse\_Figures} folder and added to this
project's \texttt{graphicspath} --- do this before finalising this section,
and include the comparison figure here. The back-propagation-of-pure-noise
variant (\cref{sec:hyp3-signbit}) has real, completed gprMax output in
\texttt{noise\_study/backprop/\{purenoise\_peaknorm,purenoise\_signbit\}/}
but its focus-frame snapshots have not yet been loaded, plotted, or
compared quantitatively --- this is the single most direct test of whether
sign-bit time-reversal actually suppresses the false-positive risk
identified above for back-propagation specifically, and is currently
missing.}

\section{Sign-Bit Time-Reversal for Noise-Robust Back-Propagation}
\label{sec:hyp3-signbit}

Back-propagation migration is structurally different from Kirchhoff and
Gazdag: it is not a post-processing step on an already-recorded image, but
requires re-injecting the (time-reversed) recorded data as a source into a
new forward simulation. This makes it vulnerable to noise in a way the
other two methods are not: spatial focusing during back-propagation is
governed almost entirely by \emph{phase} (the zero-crossings of the
time-reversed wavefield), not by amplitude, yet the default excitation
scheme injects the \emph{peak-normalised} time-reversed wavefield
$u(x,\tau)$. A large-amplitude noise spike anywhere in the data is peak-normalised
along with everything else, so it is injected with the same outsized
amplitude it had in the noisy record --- letting it act as its own
competing point source during back-propagation, interfering at the true
source locations instead of being suppressed by destructive interference.

\emph{Sign-bit time-reversal} fixes this by injecting only the
\emph{sign} of the time-reversed wavefield instead of its peak-normalised
value,
\begin{equation}
  u_{\mathrm{sign}}(x,\tau) = \operatorname{sign}\bigl(u(x,\tau)\bigr) ,
  \label{eq:signbit}
\end{equation}
implemented in \texttt{write\_backprop\_files(\ldots, sign\_bit=True)}
(\texttt{helper\_functions/migration.py}). This keeps every zero-crossing
and phase trend of the GPR wavelet completely intact, since the sign of a
signal carries its full phase information, while squashing every noise
spike down to the same $\pm1$ amplitude as the coherent signal --- stripping
noise of the outsized amplitude that would otherwise let it dominate the
back-propagated wavefield. The clean-data experiments of \cref{ch:hyp1,ch:hyp2}
use the default peak-normalised excitation throughout, since they have no
noise to suppress; every noisy back-propagation result in this chapter uses
sign-bit excitation instead.

\section{Noisy Lateral Movement}
\label{sec:hyp3-lateral}
\label{sec:tl-noise}

The lateral displacement sweep of \cref{sec:hyp1-lateral} is repeated on
B-scans contaminated with the synthetic Laplace noise of
\cref{fig:tl-bscans}c, for the two analytically-defined migration algorithms
(Kirchhoff and Gazdag; back-propagation was not re-run on this particular
noisy dataset due to its computational cost --- contrast with the vertical,
diagonal, and fluid-front cases below, where it was).
\Cref{fig:tl-noisy-kirchhoff,fig:tl-noisy-gazdag} repeat the migration and
differencing analysis of \cref{sec:hyp1-lateral} on this noisy data.

\begin{figure}[htbp]
  \centering
  \begin{subfigure}[t]{0.48\textwidth}
    \includegraphics[width=\linewidth]{TL_024_Kirchhoff_Migration_-_All_8_Noisy_Datasets____f_c15_GHz____a.png}
    \caption{All 8 noisy scenarios}
  \end{subfigure}
  \hfill
  \begin{subfigure}[t]{0.48\textwidth}
    \includegraphics[width=\linewidth]{TL_025_Kirchhoff_Migration_-_Noisy_zoomed____f_c15_GHz____aperture4.png}
    \caption{Zoomed}
  \end{subfigure}

  \vspace{0.6em}
  \begin{subfigure}[t]{0.48\textwidth}
    \includegraphics[width=\linewidth]{TL_026_Kirchhoff_Migration_-_Noisy_TimeLapse_Differences_migrated_-.png}
    \caption{Time-lapse difference}
  \end{subfigure}
  \hfill
  \begin{subfigure}[t]{0.48\textwidth}
    \includegraphics[width=\linewidth]{TL_027_Kirchhoff_Migration_-_Noisy_TimeLapse_Differences_zoomed.png}
    \caption{Difference, zoomed}
  \end{subfigure}
  \caption{Kirchhoff migration of the \emph{noisy} lateral time-lapse data
    ($f_c=\SI{15}{\GHz}$, aperture $=40$): (a,b) migrated image for all
    scenarios and zoomed; (c,d) time-lapse difference and zoomed.}
  \label{fig:tl-noisy-kirchhoff}
\end{figure}

\begin{figure}[htbp]
  \centering
  \begin{subfigure}[t]{0.48\textwidth}
    \includegraphics[width=\linewidth]{TL_028_Gazdag_Phase-Shift_Migration_-_All_8_Noisy_Datasets____f_c15.png}
    \caption{All 8 noisy scenarios}
  \end{subfigure}
  \hfill
  \begin{subfigure}[t]{0.48\textwidth}
    \includegraphics[width=\linewidth]{TL_029_Gazdag_Phase-Shift_Migration_-_Noisy_zoomed____f_c15_GHz.png}
    \caption{Zoomed}
  \end{subfigure}

  \vspace{0.6em}
  \begin{subfigure}[t]{0.48\textwidth}
    \includegraphics[width=\linewidth]{TL_030_Gazdag_Migration_-_Noisy_TimeLapse_Differences_migrated_-_mi.png}
    \caption{Time-lapse difference}
  \end{subfigure}
  \hfill
  \begin{subfigure}[t]{0.48\textwidth}
    \includegraphics[width=\linewidth]{TL_031_Gazdag_Migration_-_Noisy_TimeLapse_Differences_zoomed.png}
    \caption{Difference, zoomed}
  \end{subfigure}
  \caption{Gazdag phase-shift migration of the \emph{noisy} lateral
    time-lapse data ($f_c=\SI{15}{\GHz}$): (a,b) migrated image and zoomed;
    (c,d) time-lapse difference and zoomed.}
  \label{fig:tl-noisy-gazdag}
\end{figure}

\draftnote{compare \cref{fig:tl-noisy-kirchhoff,fig:tl-noisy-gazdag} against
the clean-data results of \cref{fig:tl-kirchhoff,fig:tl-gazdag} and state at
which displacement scale the additive noise first prevents the time-lapse
difference from being distinguished from background clutter.}

The same noisy lateral dataset is also carried through the phase-plane
estimator of \cref{sec:meth-phaseplane} (\cref{fig:tlp-noise}), for the two
migration methods available on this noisy dataset.

\begin{figure}[htbp]
  \centering
  \begin{subfigure}[t]{0.48\textwidth}
    \includegraphics[width=\linewidth]{TLP_008_NoisyTimeLapse__Phase-plane_shift_estimation__Kirchhoff____B.png}
    \caption{Kirchhoff}
  \end{subfigure}
  \hfill
  \begin{subfigure}[t]{0.48\textwidth}
    \includegraphics[width=\linewidth]{TLP_009_NoisyTimeLapse__Phase-plane_shift_estimation__Gazdag____Base.png}
    \caption{Gazdag}
  \end{subfigure}
  \caption{2D phase-plane shift estimation applied to the \emph{noisy}
    lateral-displacement dataset.}
  \label{fig:tlp-noise}
\end{figure}

\draftnote{quantify, for both methods, how much the additive Laplace noise
degrades the estimated $\Dx$ relative to the clean result of
\cref{fig:tlp-horiz-validation}, and relate this to the amplitude-weighting
argument of \cref{sec:th-mask-weight} --- this figure pair is currently the
\emph{only} direct evidence available anywhere in this thesis for the
phase-plane estimator's noise resilience; everything in
\cref{sec:hyp3-vertical,sec:hyp3-diagonal,sec:hyp3-fluidflow} below is, for
now, amplitude-level only (see \cref{sec:hyp3-gaps}).}

\section{Noisy Vertical Movement}
\label{sec:hyp3-vertical}

The vertical displacement sweep of \cref{sec:hyp1-vertical} is repeated
under the same Laplace noise, for all three migration algorithms --- unlike
the lateral case, back-propagation \emph{is} re-run here, using sign-bit
time-reversal (\cref{sec:hyp3-signbit}) on the noisy traces.
\Cref{fig:vtl-noisy-kirchhoff,fig:vtl-noisy-gazdag,fig:vtl-noisy-backprop}
repeat the migration and differencing analysis of \cref{sec:hyp1-vertical}
on this noisy data.

\begin{figure}[htbp]
  \centering
  \begin{subfigure}[t]{0.48\textwidth}
    \includegraphics[width=\linewidth]{VTL_006_B-scan_Frequency_Spectra_--_Clean_vs_Noisy____f_c15_GHz.png}
    \caption{Clean vs.\ noisy spectra}
  \end{subfigure}
  \hfill
  \begin{subfigure}[t]{0.48\textwidth}
    \includegraphics[width=\linewidth]{VTL_021_Kirchhoff_Migration_-_All_7_Noisy_Datasets____f_c15_GHz____a.png}
    \caption{Kirchhoff, all 7 noisy scenarios}
  \end{subfigure}

  \vspace{0.6em}
  \begin{subfigure}[t]{0.48\textwidth}
    \includegraphics[width=\linewidth]{VTL_026_Kirchhoff_Migration_Noisy_zoomed____f_c15_GHz____aperture40.png}
    \caption{Kirchhoff, zoomed}
  \end{subfigure}
  \hfill
  \begin{subfigure}[t]{0.48\textwidth}
    \includegraphics[width=\linewidth]{VTL_029_Gazdag_Phase-Shift_Migration_-_All_7_Noisy_Datasets____f_c15.png}
    \caption{Gazdag, all 7 noisy scenarios}
  \end{subfigure}
  \caption{Noisy vertical time-lapse data: (a) clean-vs-noisy frequency
    spectra; (b,c) Kirchhoff migration; (d) Gazdag migration.}
  \label{fig:vtl-noisy-kirchhoff}
\end{figure}

\begin{figure}[htbp]
  \centering
  \begin{subfigure}[t]{0.48\textwidth}
    \includegraphics[width=\linewidth]{VTL_030_Gazdag_Phase-Shift_Migration_Noisy_zoomed____f_c15_GHz.png}
    \caption{Gazdag, zoomed}
  \end{subfigure}
  \hfill
  \begin{subfigure}[t]{0.48\textwidth}
    \includegraphics[width=\linewidth]{VTL_034_Back-Propagation_E_Noisy_--_All_7_Datasets____focus_at_1906.png}
    \caption{Back-prop $\lVert\mathbf{E}\rVert$ (sign-bit), all scenarios}
  \end{subfigure}

  \vspace{0.6em}
  \begin{subfigure}[t]{0.48\textwidth}
    \includegraphics[width=\linewidth]{VTL_036_Back-Propagation_Ez_Noisy_--_All_7_Datasets____focus_at_1906.png}
    \caption{Back-prop $E_z$ (sign-bit), all scenarios}
  \end{subfigure}
  \hfill
  \begin{subfigure}[t]{0.48\textwidth}
    \includegraphics[width=\linewidth]{VTL_038_Back-Propagation_Noisy_--_TimeLapse_Differences_Ez____focus.png}
    \caption{Back-prop $E_z$ time-lapse difference}
  \end{subfigure}
  \caption{Gazdag (a) and sign-bit back-propagation (b--d) migration of the
    \emph{noisy} vertical time-lapse data, focused at $t=\SI{1906}{\ns}$.}
  \label{fig:vtl-noisy-gazdag}
\end{figure}

\begin{figure}[htbp]
  \centering
  \includegraphics[width=0.6\linewidth]{VTL_039_Back-Propagation_Noisy_--_TimeLapse_Differences_Ez_zoomed.png}
  \caption{Sign-bit back-propagation $E_z$ time-lapse difference for the
    noisy vertical dataset, zoomed around the scatterer depth.}
  \label{fig:vtl-noisy-backprop}
\end{figure}

\draftnote{\textbf{Gap:} this section currently only carries the analysis to
the amplitude/migrated-image level, exactly mirroring
\cref{sec:hyp1-vertical}'s clean-data detectability summary --- the 2D WLS
phase-plane fit has not yet been applied to this noisy vertical dataset.
Doing so, and comparing the result against both the clean vertical fit of
\cref{sec:tlp-vertical} and the noisy lateral fit of \cref{fig:tlp-noise},
is necessary to claim Hypothesis 3 for the vertical direction specifically.}

\section{Noisy Diagonal Movement}
\label{sec:hyp3-diagonal}

The diagonal displacement sweep of \cref{sec:hyp1-diagonal} is likewise
repeated under Laplace noise, for all three migration algorithms, again
using sign-bit time-reversal for the noisy back-propagation run.
\Cref{fig:dtl-noisy} shows the noisy Kirchhoff and Gazdag migrations and
time-lapse differences, the sign-bit time-reversed excitation itself, and
the noisy sign-bit back-propagation result.

\begin{figure}[htbp]
  \centering
  \begin{subfigure}[t]{0.48\textwidth}
    \includegraphics[width=\linewidth]{DTL_025_Kirchhoff_Migration_-_All_6_Noisy_Datasets____f_c15_GHz____a.png}
    \caption{Kirchhoff, all noisy scenarios}
  \end{subfigure}
  \hfill
  \begin{subfigure}[t]{0.48\textwidth}
    \includegraphics[width=\linewidth]{DTL_027_Kirchhoff_Migration_-_Noisy_TimeLapse_Differences_migrated_-.png}
    \caption{Kirchhoff, time-lapse difference}
  \end{subfigure}

  \vspace{0.6em}
  \begin{subfigure}[t]{0.48\textwidth}
    \includegraphics[width=\linewidth]{DTL_029_Gazdag_Phase-Shift_Migration_-_All_6_Noisy_Datasets____f_c15.png}
    \caption{Gazdag, all noisy scenarios}
  \end{subfigure}
  \hfill
  \begin{subfigure}[t]{0.48\textwidth}
    \includegraphics[width=\linewidth]{DTL_031_Gazdag_Migration_-_Noisy_TimeLapse_Differences_migrated_-_mi.png}
    \caption{Gazdag, time-lapse difference}
  \end{subfigure}

  \vspace{0.6em}
  \begin{subfigure}[t]{0.48\textwidth}
    \includegraphics[width=\linewidth]{DTL_033_Sign-Bit_Time-Reversed_Excitation_--_B-scans_and_Spectra_Noi.png}
    \caption{Sign-bit excitation B-scans \& spectra}
  \end{subfigure}
  \hfill
  \begin{subfigure}[t]{0.48\textwidth}
    \includegraphics[width=\linewidth]{DTL_038_Back-Propagation_Noisy_--_TimeLapse_Differences_Ez____focus.png}
    \caption{Sign-bit back-prop, time-lapse difference}
  \end{subfigure}
  \caption{Noisy diagonal time-lapse data: Kirchhoff (a,b) and Gazdag (c,d)
    migration and differencing; (e) the sign-bit time-reversed
    back-propagation excitation; (f) the resulting noisy back-propagation
    time-lapse difference.}
  \label{fig:dtl-noisy}
\end{figure}

\draftnote{as for the vertical case, the diagonal WLS phase-plane fit has
not yet been run on this noisy dataset --- this section is amplitude-level
only. Read off the diagonal-PSF detectability floor under noise (mirroring
\cref{sec:dtl-detectability}) and state whether it shifts relative to the
clean-data floor.}

\section{Noisy Fluid Front}
\label{sec:hyp3-fluidflow}

Finally, the fluid-front scenario of \cref{ch:hyp2} is repeated under the
same Laplace noise model, again with sign-bit time-reversal for the noisy
back-propagation run. \Cref{fig:ff-noisy} shows the noisy Kirchhoff and
Gazdag migrations and time-lapse differences, and the noisy sign-bit
back-propagation result.

\begin{figure}[htbp]
  \centering
  \begin{subfigure}[t]{0.48\textwidth}
    \includegraphics[width=\linewidth]{FF_017_Kirchhoff_Migration_-_All_8_Noisy_Datasets____f_c15_GHz____a.png}
    \caption{Kirchhoff, all noisy scenarios}
  \end{subfigure}
  \hfill
  \begin{subfigure}[t]{0.48\textwidth}
    \includegraphics[width=\linewidth]{FF_019_Kirchhoff_Migration_-_Noisy_TimeLapse_Differences_migrated_-.png}
    \caption{Kirchhoff, time-lapse difference}
  \end{subfigure}

  \vspace{0.6em}
  \begin{subfigure}[t]{0.48\textwidth}
    \includegraphics[width=\linewidth]{FF_021_Gazdag_Phase-Shift_Migration_-_All_8_Noisy_Datasets____f_c15.png}
    \caption{Gazdag, all noisy scenarios}
  \end{subfigure}
  \hfill
  \begin{subfigure}[t]{0.48\textwidth}
    \includegraphics[width=\linewidth]{FF_023_Gazdag_Migration_-_Noisy_TimeLapse_Differences_migrated_-_mi.png}
    \caption{Gazdag, time-lapse difference}
  \end{subfigure}

  \vspace{0.6em}
  \begin{subfigure}[t]{0.48\textwidth}
    \includegraphics[width=\linewidth]{FF_034_Back-Propagation_E_Noisy_--_All_Datasets____focus_at_1906_ns.png}
    \caption{Sign-bit back-prop $\lVert\mathbf{E}\rVert$}
  \end{subfigure}
  \hfill
  \begin{subfigure}[t]{0.48\textwidth}
    \includegraphics[width=\linewidth]{FF_038_Back-Propagation_Noisy_--_TimeLapse_Differences_Ez____focus.png}
    \caption{Sign-bit back-prop, time-lapse difference}
  \end{subfigure}
  \caption{Noisy fluid-front data: Kirchhoff (a,b) and Gazdag (c,d)
    migration and differencing; sign-bit back-propagation field magnitude
    (e) and time-lapse difference (f).}
  \label{fig:ff-noisy}
\end{figure}

\draftnote{the noisy fluid-front phase-plane fit (the noisy analogue of
\cref{fig:ff-phaseplane}) has not yet been run --- this is the most
field-relevant gap in this chapter, since it is the only case combining
noise robustness \emph{and} the geometry/material-change decoupling that
the unifying hypothesis explicitly claims together.}

\section{Outstanding Quantitative Work}
\label{sec:hyp3-gaps}

This chapter currently demonstrates noise robustness convincingly only at
the level of migrated \emph{images} (do they still show a detectable,
unambiguous focus under noise?), for all four cases, and demonstrates the
phase-plane estimator's noise robustness quantitatively for only one case
(noisy lateral movement, \cref{fig:tlp-noise}). To fully support Hypothesis
3, the following quantitative work, proposed but not yet run (see
\texttt{Markdown/Robustness\_Experiments.md}), is still needed:
\begin{enumerate}
  \item Apply the 2D WLS phase-plane fit to the noisy vertical, diagonal,
    and fluid-front datasets of \cref{sec:hyp3-vertical,sec:hyp3-diagonal,sec:hyp3-fluidflow},
    exactly as already done for the lateral case.
  \item Quantitatively compare sign-bit versus peak-normalised
    back-propagation on pure noise (\cref{sec:hyp3-purenoise}), using the
    completed gprMax output already sitting in
    \texttt{noise\_study/backprop/}.
  \item An orthogonality test isolating a pure geometric shift from a pure
    material-change phase rotation under noise.
  \item A controlled SNR sweep from $+30\,\si{dB}$ to $-10\,\si{dB}$
    comparing ordinary least squares against the WLS estimator.
  \item A $20\times20$ grid sweep of true $(\Dt,\Dtheta)$ pairs, reported as
    a 2D inversion-error heatmap.
\end{enumerate}
